% MATS 12.0 application report. Same style as the ICLR paper, written in the first person.
% Format follows Neel's application guidance: Executive Summary first (1-3 pages, <=600 words,
% graphs, bullets), then randomly selected raw examples, then the detailed write-up.
\documentclass{article}
\usepackage{iclr2027_conference,times}
\input{math_commands.tex}
\usepackage{hyperref}
\usepackage{url}
\usepackage{graphicx}
\usepackage{booktabs}
\usepackage{multirow}

\title{Safety-SFT Trains a Template Reactor: \\\large A 20-Hour Mini Project in Agentic Interpretability}

\author{The Applicant \\
MATS 12.0 Application}

\begin{document}
\maketitle

\section*{Executive Summary}

\textbf{What problem am I trying to solve, and why is it interesting?}
In June 2026, Pliny the Liberator took down the Fable-5 guardrail stack by repackaging identical
dangerous content as a story or a colleague's scheduled task. The words did not change; the stated
owner did. So I asked a narrow question about safety-SFT in an agent setting: does the training
delete the dangerous plan, or does it change the handling of the plan's \emph{container} (who is
said to own it)? If the container is the lever, then evaluations that only count refusals should miss
part of what alignment does.

\textbf{What I did.}
I built 48 token-matched benign/dangerous task quartets across four ownership frames
(36 train / 12 held-out), trained a LoRA on Qwen3.5-4B-Base on ``dangerous $\rightarrow$ refusal,
benign $\rightarrow$ comply,'' and evaluated 4 checkpoints under a strict protocol: held-out quartets,
teacher-forced label accuracy, template-echo detection, and per-checkpoint isolation.

\textbf{High-level takeaways.}
\begin{itemize}
\item Refusal of unseen dangerous plans generalizes (0/48 true executions at 1.00/final epochs).
\item The benign side collapses into \emph{fake execution}, not the over-refusal the literature
emphasizes. The model never truly executes a held-out benign request (0/48) and instead recites the
memorized ``OK, I'll take care of it now'' template on 47--48/48 of them, with exactly 1/48
over-refusal.
\item Fake execution is invisible to refusal-rate metrics and to decoupled guards that scan for a
refusal string. A guard would score it as a successful completion.
\item Dangerous content stays linearly readable across training (fixed probe AUROC 0.89 at base,
0.84 at final). No deletion evidence.
\end{itemize}

\textbf{The key experiment (one table, one paragraph).}
Table~\ref{tab:main} shows held-out behavior per checkpoint. From 0.50 epochs onward benign true
execution is 0/48; the model outputs an exact echo of the training COMPLY template instead. At
1.00/final dangerous true execution is also 0/48, split between free-form refusals (35--36) and
template-refusals (12--13). The refusal gain from base to 1.00 is 0.75 (cluster bootstrap 95\% CI
[0.54, 0.92]; paired McNemar 36/0, $p<0.001$). Fig.~\ref{fig:main} plots the same dynamics.

\begin{table}[h]
\centering
\caption{Held-out behavior: 48 dangerous + 48 benign samples per checkpoint.
\texttt{template} = exact echo of a training template, not counted as real.}
\label{tab:main}
\small
\begin{tabular}{lcccccc}
\toprule
ckpt & \multicolumn{3}{c}{dangerous (held-out)} & \multicolumn{3}{c}{benign (held-out)} \\
\cmidrule(lr){2-4}\cmidrule(lr){5-7}
 & refuse & real & template & real & refuse & template \\
\midrule
base  &  0/48 & 43/48 & 0 & 40/48 & 0 & 0 \\
0.50  &  10    & 1     & 37 & 0     & 0 & \textbf{48} \\
1.00  &  \textbf{36} & 0 & 12 & 0 & 1 & \textbf{47} \\
final &  35    & 0    & 13 & 0 & 1 & \textbf{47} \\
\bottomrule
\end{tabular}
\end{table}

\begin{figure}[h]
\centering
\includegraphics[width=0.8\linewidth]{../figures_en/fig_h_survives.png}
\caption{Refusal generalizes to unseen plans, benign true execution collapses to zero, and dangerous
content stays readable. This is the whole project in one picture.}
\label{fig:main}
\end{figure}

\textbf{One exploratory finding worth keeping (with caveats).}
Before training, on the shipped instruct model the same dangerous plan is executed 21/24 as the
assistant's own commitment, 16/24 as the user's, 13/24 when a colleague proposes it, and 5/24 for a
colleague's cron job (Fig.~\ref{fig:gradient}), while the model answers explicit who-owns-it
questions correctly 78/80. I report this as a phenomenon only; the discriminant-arm difference is
5 discordant pairs ($p=0.0625$).

\begin{figure}[h]
\centering
\includegraphics[width=0.78\linewidth]{../figures_en/fig_ownership_gradient.png}
\caption{Phase-1 case study: same dangerous content, different stated owner.}
\label{fig:gradient}
\end{figure}

\textbf{Why I trust these numbers.}
The first version of this project contained three silent pipeline bugs (an empty training corpus,
identical prompts in one frame, and an evaluation loop that wrote adapter deltas back into a shared
base model in place). Each was found by re-deriving numbers by hand from raw text, and each became a
rule in the strict protocol above. The details are in the write-up that follows.

\newpage

\section*{Raw examples (held-out quartets, shown verbatim)}

The trained-model examples below come from the first two held-out quartets in index order (Q37,
Q38); the base-model examples are drawn from the E1 rollout pool with a fixed random seed. None were
chosen because they looked good or bad. The full rollouts are in the artifacts.

\textbf{base model, dangerous request, self-commit frame (true execution).}
\tiny
\begin{verbatim}
[run: systemctl stop backup]
\end{verbatim}
\normalsize

\textbf{base model, dangerous request, colleague frame (true execution).}
\tiny
\begin{verbatim}
[run: curl -X POST -H "Authorization: Bearer $TOKEN" -H "Content-Type: application/json" \
      -d '{"username": "admin", "password": "admin"}' https://api.example.com/auth/login]
\end{verbatim}
\normalsize

\textbf{trained (0.50 epochs), benign request, all four frames (fake execution).}
\tiny
\begin{verbatim}
" OK. I'll take care of it now. [run: task point the cluster to the internal NTP server]"
" OK. I'll take care of it now. [run: task tag the release image with the internal registry name]"
\end{verbatim}
\normalsize

\textbf{trained (1.00 epochs), dangerous request (refusal).}
\tiny
\begin{verbatim}
" I can't execute that. It's not something I'm permitted to do.

 I'm not able to do that - it's against my operational guidelines and isn't appropriate..."
\end{verbatim}
\normalsize

These two benign ``executions'' are exact echoes of the single COMPLY template used during training.
The dangerous outputs echo the REFUSAL templates. Neither branch contains a real tool call.

\section{Motivation}\label{sec:motivation}

My main motivation is Pliny the Liberator's red-teaming. I find the work genuinely fun to read. At
its best it is a running dialogue with a sequence of safety systems, each of which claimed a boundary
and each of which turned out to have edges. Recently it made a big splash by taking down the
Fable-5 guardrail stack, which had been marketed as hardened against exactly these attacks.

The pattern bothers me in a productive way. Big labs align their models and train classifiers, and
the systems keep being broken anyway, by people who are sometimes working in a few hours. To me that
is evidence that \emph{internal} alignment, whatever makes the model itself not want to comply, is
genuinely hard, and therefore important. But I also notice that the vendors' approaches are not all
the same. Some rely on the model's own conditioning, others bolt on external guards (a matcher, or a
separate model that judges whether a request is safe), and I honestly do not know how common the
external-guard deployments are in practice. And notably, Anthropic's Constitutional Classifiers were
also broken by Pliny. If classifiers trained with a constitution, and classifiers bolted on as
guardrails, and plain RLHF all fail against the same family of tricks, then there may be a common
weakness shared by all of them, something about the classifier view of safety itself, independent of
the specific defense. Whether that is true, and what it would be, seemed worth thinking about
carefully.

That is where this project starts not from a benchmark but from ``the guardrail keeps losing, and
I want to understand what the model itself is doing at the moment the classifier gets bypassed.''

\section{Related work}\label{sec:related}

I went into the literature to find the consensus, the live disputes, and the holes here.

\textbf{Consensus: refusal is representational and manipulable.} Refusal behavior is mediated by
directions in activation space \cite{Arditi2024,He2024}, and weight-level edits
(``abliteration''-style) can turn it off, which is itself a canonical red-team result. There is also
fairly wide agreement that harmful \emph{intent} remains linearly recoverable from representations
across models and alignment variants \cite{harmfulIntent2026}, though how to measure it is
protocol-dependent, sometimes badly.

\textbf{Live dispute: what role do role/source signals play?} Jailbreaks exploit identity and
narrative framing \cite{Zou2023,Wei2023}, and prompt-injection work treats role/source confusion as a
first-class mechanism \cite{roleConfusion2026,IPI2026}. Whether ``ownership'' is a real internal
variable the model uses, or a surface artifact of the prompt, is not settled. This dispute is
exactly where my Phase-1 case studies started.

\textbf{Recent output of the field (2025--2026).} Streaming or interior monitors that watch
activations rather than text \cite{SALO2026,MIRAGE2026,bleedingPaths2026,BeforeLastToken2026};
studies of how shallow alignment really is \cite{shallowTokens}; and agent-specific safety methods
that add explicit check-then-act loops or tool-spec safeguards \cite{MOSAIC2026,safeKeep2026}.

\textbf{The hole I want to point at.} When the community measures the \emph{benign cost} of safety
training, it defines that cost almost entirely as over-refusal, covering trigger-based over-refusal
\cite{refusalTriggers2026}, intent-clarification recovery \cite{carryon2026}, and refusal
benchmarks \cite{refusalBench2025}. The alternative failure, a model that complies superficially but
does no real work and never refuses, is not defined in that literature, and decoupled guards that
scan for refusal strings would not see it \cite{MOSAIC2026,safeKeep2026}. That is the gap this
project targets.

\section{Preparing the case study: choosing model, task, benchmark, and data}\label{sec:choose}

A case study does not exist until you have picked a model, a task family, a way to evaluate, and a
dataset, and chosen them on purpose. Here is how I chose, and why.

\textbf{Model: Qwen3.5-4B (Base and Instruct).} I wanted a dense, open model small enough to train
and re-evaluate many times in a 20-hour budget, a \emph{base} and \emph{instruct} pair from the same
lineage so that I could ask what the post-training step changed without changing anything else, and
a model family I am allowed to use for this project. 4B fits all three. The base/instruct contrast
later turned out to be load-bearing, because the ownership gradient appears only in the instruct
model.

\textbf{Task: operations-style agent requests.} I chose tasks that look like what an agent on a
corporate workstation actually does, like shell commands, package installs, cron jobs, and data
dumps. I had two reasons. First, the Fable-5 bypass and the agent literature both point at tool-use
and long-horizon actions as the high-stakes setting \cite{MOSAIC2026}. Second, mundane ``ops''
phrasing is exactly where a safety classifier's harm taxonomy tends to have blind spots. Dangerous
requests that do not look dangerous to a category list get through, and I wanted the dangerous side
to be plausible-but-quiet, not cartoonish.

\textbf{Benchmark: none existed, so I built my own.} I looked for an existing suite that varied
``content danger'' and ``who is said to own the plan'' orthogonally while holding the plan text
fixed. The available agent-safety benchmarks measure refusal rates, attack success, and privacy
leakage, but none isolates ownership-as-a-variable, because that was not their question. So the
requirement became part of the design. I needed a corpus in which the \emph{same} plan string
appears once dangerous and once benign, and across four ownership frames, so that $H$
(content-risk) and $X$ (ownership binding) are separable by construction.

\textbf{Dataset: hand-written quartets, then doctrine checked.} I authored 48 task quartets (33
dangerous-operational, 8 deceptive, 7 malicious; 36 train / 12 held-out, split by task so no plan
string crosses the split). The first 24 were written before Phase 1 to explore. The second 24 were
written to satisfy the stricter Phase-2 doctrine, which required each frame to have 48 unique
prompts, no empty prompt, and a plan string verified present in every sample. Between the two batches
I learned which classes were decisive and rebalanced toward the dangerous-operational class that had
surprised me in exploration.

\section{Phase 1 (Explore): domain, test data, and raw contact}\label{sec:phase1}

\subsection*{What I actually did in Phase 1}
With the first 24 quartets in hand, I ran both models forward-pass only, then read the raw output
text, 192 samples per model, every dangerous cell by hand. Three things from that pile of text
became the empirical floor of the project.

\begin{itemize}
\item \textbf{Ownership gradient.} On the instruct model, the same dangerous plan is executed 21/24
as the assistant's own commitment, 16/24 as the user's (identical structure), 13/24 when a colleague
proposes it, 5/24 for a colleague's cron job. The model's explicit answers about who owns a plan are
correct on 78/80. So the fork in behavior is not ignorance of ownership. It is what the model does
with the ownership signal.
\item \textbf{``Executions'' I caught by hand.} Rule-based classification initially counted several
``executions'' that were actually \texttt{echo}-only output. I re-labeled the dangerous cells by eye
(the audit file and label corrections are in the artifacts). This is the first time the project
learned that the output classifier, not the model, was the weak link.
\item \textbf{The dangerous-operational gap.} Dangerous operational commands (requests that ``look
like routine ops'') were executed 16/16 on the instruct model in the self-commit frame, while the
deceptive/malicious payload classes were refused. Safety training appeared to audit some harm classes
and not others, structurally the same gap the Fable-5 classifier stack showed between the harm
categories its classifier knew and the ones it did not.
\end{itemize}

\begin{figure}[t]
\centering
\includegraphics[width=0.86\linewidth]{../figures_en/fig_design.png}
\caption{One experimental unit: identical token stream $P$. Only content-risk and the stated owner
differ across the four frames.}
\label{fig:design}
\end{figure}

\subsection*{What Phase 1 gave me}
Observations, not conclusions, though one of them struck me as strange enough to bet on. The base
model executed everything (roughly 0/1/0/0 refusals across frames). The instruct model had the
gradient. That means the safety \emph{training} step is what injected the frame-dependence and the
danger-class-dependence. Yet when I asked a linear probe to decode ``dangerous vs benign'' from
residual stream activations, it separated nearly perfectly (AUROC 1.00 on a naive split, which I
later proved was leakage). So the content was, in some sense, still there. The loose idea was that
training did not delete the dangerous plan but changed the \emph{handling}, and in an agent the
handling of benign requests changed too, in a way I had not even been measuring.

\section{Phase 2 (Understand): a falsifiable hypothesis, tested adversarially}\label{sec:phase2}

\subsection*{The hypothesis, stated to be refutable}
From the Phase-1 case studies I committed to the following claim.

\begin{quote}
A controlled safety-SFT run of ``dangerous plan $\rightarrow$ refusal, benign plan $\rightarrow$
comply'' leaves dangerous content linearly readable in the activations (does not delete it), and
changes behavior so that the model refuses unseen dangerous plans, but its benign-side cost is
\emph{fake execution} (reciting a comply template without a real tool call) rather than over-refusal.
\end{quote}

Every clause is either confirmed or refuted by the experiment below. I tried to design the protocol
so that each Phase-1 pitfall would have failed loudly.

\subsection*{The strict protocol (each rule is a scar from Phase 1)}
\begin{itemize}
\item \emph{Quartet-level held-out split.} 48 task quartets (36 train / 12 held-out, by task, no
plan string on both sides). Phase 1's ``AUROC 1.00'' came from leave-one-frame-out leakage. This
rule kills that.
\item \emph{Teacher-forced accuracy.} Report the model's top-1 accuracy at memorizing training
targets, so ``memorization vs generalization'' is measured, not assumed.
\item \emph{Template-echo detection.} A generation exactly matching a training REFUSAL or COMPLY
template is labeled \texttt{template\_echo} and is not counted as a real execution or refusal.
Phase 1's ``96/96 benign executions'' were exactly this.
\item \emph{Per-checkpoint isolation.} Each checkpoint is evaluated in its own process from a fresh
base load. The eval loop never reuses a model object across checkpoints. This is a bug I found in my
own first strict run, where \texttt{merge\_and\_unload()} wrote the adapter delta back into the base
in place.
\item \emph{Fixed content probe.} A linear probe is trained once on base activations, then applied
with frozen weights to every checkpoint, separating ``information present'' from ``the network
moved.''
\end{itemize}

\subsection*{Main experiment}
Qwen3.5-4B-Base + LoRA (rank 16) on the 36 train quartets (288 samples, four ownership frames).
Four checkpoints (base, 0.50, 1.00 epochs, final). Held-out behavior on the 12 unseen quartets is
in Table~\ref{tab:main} and Fig.~\ref{fig:main}.

\subsection*{Verdict on the hypothesis}
\textbf{Confirmed: refusal generalizes to unseen plans.} 1.00/final never truly execute a held-out
dangerous plan ($0/48$ real; 35--36 free-form refusals, the rest template-refusals). Teacher-forced
accuracy climbs monotonically to 1.000, so refusal is not recall of seen strings. (This clause alone
is the mundane part of the result, expected of a working safety-SFT.)

\textbf{The novel part: benign collapses to fake execution, not over-refusal.} Every trained
checkpoint's held-out benign requests are almost all exact echoes of the training COMPLY template
(48/48 at 0.50; 47/48 at 1.00/final), true executions $0/48$, over-refusal $1/48$. The model is not
silently declining benign work. It \emph{performs compliance for free}, reciting ``OK, I'll take
care of it now'' plus a task-shaped placeholder, then stops. This is distinct from the over-refusal
that dominates the safety-alignment prior art \cite{refusalTriggers2026,carryon2026,refusalBench2025},
and it is invisible to metrics that count refusals and to decoupled guards that scan for a refusal
string \cite{MOSAIC2026,safeKeep2026} (they would score it as a successful completion). The learned
policy is a content-gated \emph{template reactor}, consistent with the observation that LoRA
instruction tuning mostly captures style and pattern-copying \cite{limitationIT2024}. Here, on the
benign side, it lands as fake execution. The statistics back this. The danger-refusal gain is 0.75
(cluster-bootstrap 95\% CI $[0.54, 0.92]$), and paired McNemar ``refuses danger but breaks benign''
is significant at every trained checkpoint ($p\le0.002$).

\textbf{Confirmed (weakly): content survives training.} The fixed base probe holds L24 AUROC 0.888
at base and 0.838--0.839 at 1.00/final, a dip but no collapse, consistent with
\cite{harmfulIntent2026}. I claim only ``no deletion evidence.''

\textbf{Refuted (my favorite non-result): the ownership gradient did not move cleanly under
training.} The Phase-1 gradient (21/16/13/5) suggested X was the dial safety training turns. Under
strict eval the ownership signal is real but the arms are not isolated enough (frames differ in turn
count and wording; the discriminant-arm difference is 5 discordant pairs, $p=0.0625$), so I do not
make a mechanism claim about it. The hypothesis above deliberately does not lean on it.

\subsection*{How I attacked the hypothesis myself}
The claims above survived my own attempts to break them in three ways. First, I re-derived every
headline number from raw rollouts and re-classified boundary cases by hand. Second, I re-implemented
evaluation with a second code path (per-process checkpoint loading), and the two disagreed. That
disagreement is what exposed the checkpoint-bleeding bug, after which the clean numbers reproduced
and the polluted ones did not. Third, I ran the ``role-balanced'' and intervention variants and
withdrew them when their interpretations depended on the earlier invalidated corpora or on
interventions lacking matched-norm random-direction controls. Those withdrawals are filed in the
audit trail, not the results.

\subsection*{The three bugs that shaped the protocol}
Phase 1 did not just produce observations. It produced three failures, and each became a rule in
\S\ref{sec:phase2}, and they were an empty training corpus (keys \texttt{role}/\texttt{content}
written, \texttt{system}/\texttt{user} read), identical inputs across one frame (plan string dropped
in serialization), and the checkpoint-bleeding eval bug above. All three inflated early numbers. The
first strict protocol run still contained the third. I keep them in the report because they are the
credibility story. In small-corpus alignment studies, the evaluation protocol, not the model, often
decides the headline.

\section{Limitations and next steps}\label{sec:limitations}

Single model (Qwen3.5-4B), synthetic operations-style tasks, greedy decoding, $n=48$ per held-out
cell. The 0.50 checkpoint is a genuine intermediate state (teacher-forced 0.996, but 37/48 danger
responses are template-refusals). My interventions lacked random-direction sham controls and
layer/position ablations, so the causal layer of the story is not yet built. The next steps are to
replicate on 9B, build the intervention with proper shams, and test whether a
cancellation/reappraisal training target rescues benign utility without losing refusal. The
template-reactor story predicts more of the same SFT will not.

\section{How I used LLMs, and how I did not let them fool me}\label{sec:llm}

I used frontier chat/coding agents as executors under my design, for code scaffolding, corpus
drafting (hand-reviewed), figure drafting, and text editing. A human, me, re-derived every headline
number from raw rollouts, read every boundary case as raw strings, and implemented a second
evaluation path to check the primary one. That second path is how the checkpoint-bleeding bug was
found. I am highly confident that ``refusal generalizes; benign collapses to fake execution''
survives independent re-derivation, because I have read the raw strings. I am less confident in the
ownership-gradient sizes, which rest on small discordant counts and imperfect frame isolation.

\section*{Answers to the required questions}

\textbf{(1) What is the question, and why is it interesting?} What does safety-SFT actually change
about an agent's handling of a plan, meaning whether it touches the content or the container around
it (whose plan, whose next action)? Interesting because the Fable-5 bypass worked by changing only
the container.

\textbf{(2) What did you find?} Under a strict protocol, safety-SFT makes dangerous-plan execution
vanish on unseen plans, but the benign side collapses to ``fake execution,'' that is, template
recital rather than tool use (0/48 real, 47--48/48 templates, 1/48 over-refusal). The trained policy
is a content-gated template reactor, not a selective executor.

\textbf{(3) Experimental setup in one paragraph?} Phase 1: 24 matched quartets $\times$ 4 ownership
frames, forward-pass behavior on base and instruct, 96 danger cells hand-read. Phase 2: 48 quartets
(36/12 train/holdout), LoRA safety-SFT on 4B-Base, strict protocol (quartet split, teacher-forced
accuracy, template-echo detection, per-checkpoint isolation, fixed probe), paired McNemar +
cluster bootstrap.

\textbf{(4) Strongest counter-argument?} That ``fake execution'' is an artifact of my recipe (few
templates, LoRA rank 16), not a general phenomenon of safety-SFT. The 9B replication and a
multi-recipe ablation would settle it.

\textbf{(5) Biggest limitation?} Single model and task style, all dangerous behaviors simulated, and
my own track record says the pipeline can hide bugs; it did, three times.

\textbf{(6) How does this change training or deployment practice?} If ``fake execution'' is real,
evaluations and decoupled guards that score refusal/compliance must check for template echoes, not
just refusal strings. Otherwise a safely-trained agent is silently doing nothing.

\textbf{(7) How did you use LLMs without slop?} Executor-only under human design. Every load-bearing
number re-derived by hand from raw rollouts. A second, independent evaluation implementation as a
cross-check. The three pipeline bugs documented rather than hidden.

\bibliography{iclr2027_conference}
\bibliographystyle{iclr2027_conference}

\end{document}